\documentclass[a4paper, 17pt]{extarticle}
\usepackage{siunitx}
\usepackage{amsmath}
\usepackage{microtype}
\usepackage{tikz}
\usepackage{standalone}
\usepackage{tikz}
\usepackage[landscape]{geometry}
\sisetup{exponent-product=\cdot}

\title{Printable cutouts}
\date{2026--02--23}
\author{S. Bodapati \and D. Jain \and S. Malhan \and A. Shokeen}

\begin{document}
\section{Problem statement}
Redshift and blueshift are phenomena wherein light from distant objects is ``stretched out''; that is, the observed wavelength increases or decreases relative to the emitted wavelength. This affects the light received on Earth, and can be defined by the following equation:
\begin{align}
    z &= \frac{v}{c\sqrt{1 - \frac{v^2}{c^2}}}
\end{align}
wherein \(z\) is the redshift (If \(z = 0\), then there is no redshift or blueshift. If \(z < 0\), it is blueshift. And if \(z > 0\), then there is redshift.), \(v\) is the recessional velocity, in \unit{\metre\per\second}, and \(c\) is the speed of light in a vacuum, approximately \qty{299 792 458}{\metre\per\second}.

This is a problem for cosmologists and spectrographers. Cosmologists need accurate images to study the universe, and whilst tools exist to remedy this, they are nonetheless very slow and inefficient. Spectrographers --- especially extraterrestrial spectrographers --- rely on spectral lines to determine the elemental composition of planets' atmospheres, stars' atmospheres, and more. Our product aims to enhance our interpretation of redshift (Not stop it, because that would require stopping the expansion of the universe, which is beyond our budget.\textsuperscript{[citation needed]}) by accounting for the redshift digitally.

\section{Hypothesis}
If there is an increase in the distance to the object in question, then there shall be an increase in its recessional velocity, because the Hubble-Lemaître law states, mathematically:
\begin{align}
    v &= H_0 \cdot d
\end{align}
wherein \(H_0\) is the Hubble constant: \qty{2.23e-18}{\per\second}, \(d\) is the distance, and \(v\) is the recessional velocity. Thanks to this, we can say that \(v \propto d\).

\newpage

\section{Materials}
\begin{itemize}
    \item ASUS Vivobook
    \item Python version \(\pi\)
    \item Einstein's and Hubble's research
    \item VS Code
    \item Basic\TeX
    \item \LaTeX{} packages: \verb|siunitx, csquotes, microtype, amsmath|.
\end{itemize}

\newpage

{\LARGE \section{Big question}
 How can we effectively use current technology to improve cosmology and spectrography by addressing the problem of cosmological redshift and blueshift digitally using well-defined scientific principles such as special and general relativity, the Hubble-Lemaître law (\(v = H_0 \cdot d\)), and optics?}
\newpage

\section{Procedure}
We began by trying to use formulae to anti-redshift a wavelength based on the distance and observed wavelength:
\begin{align}
    v &= H_0 \cdot d \\
    z &= \frac{v + c}{c\sqrt{1 - \frac{v^2}{c^2}}} - 2\\
    \lambda_{\text{obs}} &= \lambda_{\text{emit}}(z + 1)
\end{align}
We ran into a problem, however, when we realized we could not convert directly from digital colour spaces to wavelengths, because a digital colour space such as RGB is made up of 3 channels: red, blue, and green, which do not map to a wavelength exactly.

We then decided to train an AI to remove this redshift by tricking it into thinking there was a stained glass pane on top of the image. This worked sometimes, but not all of the time, because AIs are unpredictable. This was clearly sub-optimal, and so we decided to do something else: gaussian blurring.

Gaussian blur is explained in more detail by Bodapati, but it is essentially when you blur the less impoertant bits of an image. We used this to eliminate distortion, and reduce redshift.

\end{document}
